\chapter{Sleep-Related Bruxism}
\index{Sleep-Related Bruxism}
\label{sec:Sleep-Related Bruxism}

\section{Overview}
Sleep-related bruxism is a sleep-related movement disorder characterized by repetitive jaw-muscle activity during sleep, manifesting as tooth grinding or clenching, often associated with rhythmic masticatory muscle activity (RMMA). Prevalence is 8--12\% in adults and 15--30\% in children, decreasing with age. This condition is among the most common mental health disorders worldwide. It affects people across all age groups, socioeconomic backgrounds, and geographic regions.
\section{Causes}
This condition happens because of sleep-related activation of the brainstem masticatory central pattern generator, triggered by spontaneous microarousals and possibly modulated by autonomic activation.     
\section{Risk Factors}

\begin{itemize}[noitemsep]
  \item Family history of mental illness
  \item Trauma or adverse childhood experiences
  \item Chronic stress
  \item Substance use
  \item Social isolation
  \item Underlying medical conditions
  \item Genetic predisposition and family history
  \item Advancing age
  \item Lifestyle factors (diet, exercise, smoking, alcohol)
  \item Environmental exposures
  \item Medication use
\end{itemize}
\section{Signs and Symptoms}

\subsection{Common symptoms}
\begin{itemize}[noitemsep]
  \item \item You may experience audible tooth grinding, jaw pain or fatigue, temporal headaches, tooth wear, tooth sensitivity, and temporomandibular joint pain. Pain or discomfort in the affected area    \item Pain or discomfort in the affected area
  \end{itemize}

\subsection{Emergency warning signs}
\begin{itemize}[noitemsep]
  \item Severe or worsening symptoms of sleep-related bruxism require immediate medical evaluation
\end{itemize}
\section{Progression and Stages}
The condition may progress through different stages of severity over time. Early stages often involve mild or subtle symptoms that may be overlooked. Moderate stages involve more noticeable symptoms affecting daily function. Advanced stages may involve significant impairment and complications.
\section{Complications}
\begin{itemize}[noitemsep]
  \item Disease progression leading to worsening symptoms
  \item Secondary infections or injuries
  \item Side effects of treatment
  \item Reduced quality of life
  \item Emotional and psychological impact
\end{itemize}
\section{Diagnosis}
\begin{itemize}[noitemsep]
  \item Doctors usually diagnose this based on your symptoms and a physical exam
  \item Polysomnography with masseter EMG is the gold standard.
  \item Comprehensive medical history and physical examination
  \item Laboratory tests to assess organ function
  \item Imaging studies to visualize affected structures
  \item Specialized testing depending on the affected system
  \item Consultation with relevant specialists
\end{itemize}
\section{Prevention}
\begin{itemize}[noitemsep]
  \item Maintain a healthy lifestyle with balanced diet and exercise
  \item Avoid known risk factors
  \item Attend regular medical check-ups for early detection
  \item Follow recommended screening guidelines
\end{itemize}
\section{Treatment Options}
Treatment options include occlusal splints for tooth protection, stress reduction, treatment of comorbid OSA, botulinum toxin injections for severe cases, and avoidance of exacerbating substances. Medications to manage symptoms and address underlying mechanisms Lifestyle modifications and self-care measures Physical therapy and rehabilitation Surgical intervention when indicated Regular monitoring and follow-up care Patient education and support services

\begin{itemize}[noitemsep]
  \item Medications to manage symptoms and address underlying mechanisms
  \item Lifestyle modifications and self-care measures
  \item Physical therapy and rehabilitation
  \item Surgical intervention when indicated
  \item Regular monitoring and follow-up care
\end{itemize}
\section{Living with It}
\begin{itemize}[noitemsep]
  \item Follow your treatment plan as prescribed
  \item Maintain regular follow-up with your healthcare provider
  \item Make necessary lifestyle adjustments
  \item Seek support from family, friends, or support groups
  \item Stay informed about your condition
\end{itemize}
\section{Outlook}
\begin{itemize}[noitemsep]
  \item In most cases, the outlook is manageable
  \item Severe untreated bruxism can result in tooth fracture and TMJ disorders.
\end{itemize}
